\section*{Appendix: Full O-OIS EBNF Grammar}
\addcontentsline{toc}{section}{Appendix: Full O-OIS EBNF Grammar}

This appendix gives the complete formal grammar of O-OIS in Extended Backus--Naur Form (EBNF). A reader who wants to reimplement O-OIS should start here. The main text (Section~\ref{sec:ontology}) explains the same grammar in plain language with examples. The grammar below is a specification, not a program: it says what is allowed, not how to write the code.

\subsection*{A.1 Transport layer: content-addressed handles}

Every graph and key is named by a short handle derived from its content. Two identical graphs always receive the same handle.

\begin{verbatim}
TransportHandle   ::= KeyHandle | GraphHandle | TextHandle
KeyHandle         ::= "<K:" Hex4 ">"
GraphHandle       ::= "<G:" Hex4 ">"
TextHandle        ::= "<T:" Hex4 ">"
Hex4              ::= HexDigit HexDigit HexDigit HexDigit
HexDigit          ::= "0".."9" | "a".."f"

ControlOp         ::= LoadOp | UnloadOp | NoRetainLoadOp | CondenseEvictOp
                    | ListTypeOp | DiffOp | ProveOp
LoadOp            ::= "LD" TransportHandle
UnloadOp          ::= "UNL" TransportHandle
NoRetainLoadOp    ::= "NDL" TransportHandle
CondenseEvictOp   ::= "RSP" TransportHandle
ListTypeOp        ::= "LST" TypeName
DiffOp            ::= "GAP" GraphHandle
ProveOp           ::= "PRV" GraphHandle GraphHandle

TypeName          ::= "entity" | "edge" | "facet" | "salience"
                    | "reification" | "literal"
\end{verbatim}

\textit{Plain-language note.} The handle is the first 16 bits of a SHA-256 hash of the content. At $10^4$ objects, the probability of a collision is about $10^{-5}$. Handles remain valid across program restarts. Together with parity and a checksum, they form error-correction transport: the data plus the extra information needed to detect and repair damage.

\subsection*{A.2 Core graph grammar}

\begin{verbatim}
Graph             ::= Entity* Edge* Facet* Salience* Reification*
Entity            ::= "entity" EntityId Label
                    CanonicalId? PropertiesLiteral? FacetRef* SalienceRef*
EntityId          ::= Integer (0-indexed, implicit from position)
Label             ::= UTF8String (non-empty, max 255 bytes)
CanonicalId       ::= "@" UTF8String (max 8 bytes, padded/truncated)
PropertiesLiteral ::= "{" PropEntry* "}"
PropEntry         ::= Key ":" LiteralValue
Key               ::= UTF8String
LiteralValue      ::= BoolLit | IntLit | FloatLit | StringLit
BoolLit           ::= "true" | "false"
IntLit            ::= ["-"] Digit+
FloatLit          ::= ["-"] Digit+ "." Digit+ ("e" ["-"] Digit+)?
StringLit         ::= "\"" UTF8String "\""
FacetRef          ::= "#facet" FacetId
SalienceRef       ::= "#sal" SalienceId

Edge              ::= "edge" EdgeId Source Target Relation
                    Subtype? Confidence Polarity Reified?
EdgeId            ::= Integer (0-indexed)
Source            ::= EntityId
Target            ::= EntityId | LiteralValue
Relation          ::= RelationType
RelationType      ::= "EQL" | "ISA" | "INS" | "PAR" | "CST" | "PRC" | "PRE"
                    | "LOC" | "USF" | "ANT" | "PRP" | "OTH"
Subtype           ::= MeronymySubtype | AssociativeLabel
MeronymySubtype   ::= "CIN" | "MEM" | "POR" | "STU" | "FEA" | "PLA"
AssociativeLabel  ::= "HAS_PROPERTY" | "USED_FOR" | "LOCATED_IN"
                    | "PART_OF" | "IS_A" | "OTHER"
Confidence        ::= "CERTAIN" | "LIKELY" | "PLAUSIBLE" | "UNKNOWN"
                    (maps to CNF 0..3)
Polarity          ::= "1" | "0"  (1 = asserted, 0 = negated)
Reified           ::= "reified" | epsilon

Facet             ::= "facet" FacetId DimensionId+
FacetId           ::= Integer
DimensionId       ::= Integer

Salience          ::= "salience" EntityId Source Scalar
Source            ::= "EMPIRICAL" | "DESIGNER_INTENDED" | "INFERRED" | "DEFAULT"
Scalar            ::= Float (0.0 .. 1.0)

Reification       ::= "reify" EdgeId
\end{verbatim}

\subsection*{A.3 Relation algebra}

\textit{Primary relations.} Read the table with two questions: can the relation be read backwards (symmetric)? Does it chain (transitive)?

\begin{table}[h]
\centering
\begin{tabular}{llccc}
\toprule
Symbol & Name & Symmetric & Transitive & Meaning \\
\midrule
EQL & Equivalence & Yes & Yes & Same referent \\
ISA & Is-A & No & Yes & Subsumption (A ISA B) \\
INS & Instance-Of & No & No & Instantiation (a INS A) \\
PAR & Part-Of & No & Yes & Mereological parthood \\
CST & Constitutes & No & Yes & Inverse of PAR \\
PRC & Produced-By & No & Yes & Causal / production dependency \\
PRE & Precedes & No & Yes & Temporal / logical precedence \\
LOC & Located-In & No & Yes & Spatial / topical containment \\
USF & Used-For & No & No & Teleological function \\
ANT & Antonym & Yes & No & Semantic opposition \\
PRP & Has-Property & No & No & Attribute attachment \\
OTH & Other & Yes & No & Fallback for unclassified \\
\bottomrule
\end{tabular}
\end{table}

\textit{Meronymy subtypes} refine PAR and CST. Each one says \textit{how} a part belongs to a whole:

\begin{table}[h]
\centering
\begin{tabular}{lll}
\toprule
Symbol & Expansion & Meaning \\
\midrule
CIN & Component-Integral & Essential functional component \\
MEM & Member-Collection & Interchangeable member of a collection \\
POR & Portion-Mass & Measured portion of a mass or substance \\
STU & Stuff-Object & Material constituting an object \\
FEA & Feature-Activity & Feature or activity of a process \\
PLA & Place-Area & Sub-area within a larger area \\
\bottomrule
\end{tabular}
\end{table}

\textit{Associative labels} are plain-language refinements.\\ HAS\_PROPERTY refines PRP; USED\_FOR refines USF;\\ LOCATED\_IN refines LOC; PART\_OF refines PAR;\\ IS\_A refines ISA; and OTHER is the fallback.

\subsection*{A.4 Epistemic layer}

\begin{verbatim}
EpistemicTag      ::= NegTag | ConfidenceTag | SalienceTag | FacetTag
NegTag            ::= "NEG" ("0" | "1")           -- 0 = negated, 1 = asserted
ConfidenceTag     ::= "CNF" ("0" | "1" | "2" | "3")
                    -- 0=certain, 1=likely, 2=plausible, 3=unknown
SalienceTag       ::= "SAL" Float                 -- 0.0 .. 1.0
FacetTag          ::= "FAC" FacetId
ReificationTag    ::= "REI" EdgeId                -- reify edge as node
\end{verbatim}

Rules: NEG\,=\,0 on an ISA edge means classical negation (``A is not a B''). CNF\,=\,0 means deductively certain; CNF\,=\,3 means unknown. SAL at or above 0.5 means ``include by default.'' FAC partitions the graph into overlapping points of view. REI lets an edge become a node, so that statements can be made about statements.

\subsection*{A.5 Tiered loading control}

\begin{verbatim}
LoadPolicy        ::= ResidentPolicy | PointedPolicy | CondensedPolicy
ResidentPolicy    ::= "LD" GraphHandle "resident"  -- stays in Tier-1 context
PointedPolicy     ::= "NDL" GraphHandle "pointed"
                    -- load on demand, evict after use
CondensedPolicy   ::= "RSP" GraphHandle "condensed"
                    -- summary only (stats + handles)

LoadDecision      ::= "load" GraphHandle LoadPolicy
EvictDecision     ::= "unload" GraphHandle
GapDecision       ::= "gap" GraphHandle ResidentHandle  -- diff vs resident
ProveDecision     ::= "prove" GraphHandle GraphHandle
                    -- structural equivalence
\end{verbatim}

\textit{Invariant (MSCA).} The Tier-1 resident budget is at most 640 tokens (p50k\_base). Tier-2 sections are loaded by phase and tag with \texttt{python -m msca.loader --phase N --tag T}. A non-zero exit means failure: never continue from a partial load.

\subsection*{A.6 Encoding grammar (POINTS glyph)}

\begin{verbatim}
Glyph             ::= Header Content Parity
Header            ::= Magic Version Sequence ParityInfo EntityInfo EdgeInfo
                    PayloadInfo GraphId RegionOffsets
Magic             ::= "PTS\x01" (4 bytes)
Version           ::= UInt8 (1)
Sequence          ::= SeqId SeqIndex SeqTotal
SeqId             ::= UInt32
SeqIndex          ::= UInt16
SeqTotal          ::= UInt16
ParityInfo        ::= ParityRatio ParityAlgoId
ParityRatio       ::= UInt8 (0..100)
ParityAlgoId      ::= UInt8 (1=XOR, 2=RS)
EntityInfo        ::= EntityCount EntityIndexBits ConfidenceBits
EntityCount       ::= UInt16 (0..65535)
EntityIndexBits   ::= UInt8 (1..16)
ConfidenceBits    ::= UInt8 (2 fixed)
EdgeInfo          ::= EdgeCount
EdgeCount         ::= UInt16 (0..65535)
PayloadInfo       ::= FacetCount SalienceCount LiteralPoolBytes
GraphId           ::= UInt32
RegionOffsets     ::= Offset[7] (UInt24 each)

Content           ::= EntityTable EdgeTable Salience Facets Reification
                    LabelPool LiteralPool
EntityTable       ::= EntityRecord* LabelPool
EntityRecord      ::= Flags CanonicalId? LabelOffset LabelLen PropertiesLiteral?
Flags             ::= UInt8 (bit0=label, bit1=canonical_id,
                             bit2=properties_literal)
CanonicalId       ::= Bytes[8] (UTF-8, padded)
LabelOffset       ::= UInt16
LabelLen          ::= UInt8
PropertiesLiteral ::= PropCount (KeyLen Key LiteralValueBytes)*
PropCount         ::= UInt8

EdgeTable         ::= EdgeRecord*
EdgeRecord        ::= CompactEdge | FullEdge
CompactEdge       ::= "0" Source(EIB) "0" Target(EIB)
                    RelationType(3) Polarity(1)
FullEdge          ::= "1" Source(EIB) TargetKind TargetBits RelationType(3)
                    Subtype(3) Confidence(2) Polarity(1) Reified(1)
TargetKind        ::= "0" (entity) | "1" (literal)
TargetBits        ::= max(EIB, 16)

Parity            ::= XORParity | RSParity
XORParity         ::= Bytes[parity_size] (ratio% of content)
RSParity          ::= RS(255,223) parity bytes (32 bytes per 223-byte block)
\end{verbatim}

\textit{Constants.} GLYPH\_SIZE = 8192 bytes; HEADER\_SIZE = 42 bytes; CONTENT\_SIZE = 6920 bytes; default parity ratio 15\% (XOR) or 32 parity bytes per block (RS). Maximum entities per glyph: about 460 in a simple chain, about 300 with rich edges. Larger graphs use multi-glyph chaining via sequence numbers.

\subsection*{A.7 Semantic capacity metric (formal)}

\begin{verbatim}
SC(R, G) = |sem(G)| / |R|  (bits-of-meaning per byte)

|sem(G)| = sum over p in propositions(G) of w(p) * bits(p)
  w(p) = 1 (uniform weight)
  bits(p) = 10 (conservative estimate per atomic proposition)

propositions(G) = { entity_label(e) | e in E }
                union { entity_cid(e) | e in E and cid(e) != empty }
                union { edge(e) | e in Edges }  -- 6-field tuple
                union { facet(f) | f in Facets }
                union { salience(s) | s in Salience }
\end{verbatim}

\textit{Baselines.} Compact JSON, Markdown, YAML, token stream (BPE), and random bytes. \textit{Corruption conditions:} clean round trip; RS repair at 1, 4, 8, 16, and 32 damaged bytes per 223-byte block; truncation to 10\%, 25\%, and 50\%; adversarial reordering, synonym substitution, and removal of epistemic tags.

\subsection*{A.8 Optimality scope (honest)}

O-OIS is a regulative ideal (in Kant's sense), not a provable global optimum. The claims made are: structural completeness (every known failure mode of predecessor codecs is addressed, patches B1--B4); comparative design quality (the best characterised alternative for maximising semantic capacity under fixed-size, error-correcting, human-checkable constraints, as of November 2026); and indexicality (every claim is bounded by context, date, and knowledge horizon, expiring 22 December 2026).

\subsection*{A.9 Disclosure block (per MSCA)}

\begin{verbatim}
LOADED     <G:points> <G:msca> <G:65536>
NOT LOADED (all further handles: outside this paper's scope)
\end{verbatim}

\textit{Verification.} This grammar was checked against the POINTS example set (7 entries), a capacity sweep (1 to 460 entities), and a Reed--Solomon damage test (0 to 32 damaged bytes).
